\documentclass{elsinta}
\usepackage{url}

% --- Memasukkan Data Dokumen Menggunakan Perintah Baru ---
\projectcode{EL2}
\documentname{\LiteraturesStandardsTitle}
\documenttitle{Teknik Elektro - Sistem Informasi Tugas Akhir}
\capstonetitle{Teknik Elektro - Sistem Informasi Tugas Akhir}
\shortname{  ELSINTA} 
\documentnumber{TA2526.01.099}
\revisionnumber{03}
\publicationdate{22 September 2025}

\revdatefooter{03/10/2025}
% --- Memasukkan Data Halaman Pengesahan (Halaman 2) ---
\currentpage{2} % Halaman yang sedang dicetak

% Data Tim
\ketuatimnama{Mahasiswa 1}
\ketuatimnim{12345678}

\anggotanamaI{Mahasiswa 2}
\anggotanimI{12345678}

\anggotanamaII{Mahasiswa 3}
\anggotanimII{12345678}


% Data Dosen
\pembimbingInama{Dosen 1}
\pembimbingInip{12345678 123456 1 123}

\pembimbingIInama{Dosen 2}
\pembimbingIInip{12345678 123456 1 123}


% Tanggal Pengesahan
\approvaldate{22 September 2025}

\begin{document}

% Panggil perintah untuk membuat halaman sampul
\coverpage
% Pindah ke Halaman Baru
\newpage

% Membuat Halaman Pengesahan (Halaman 2)
\approvalpage

% Lanjut ke halaman proposal
\newpage

% ======================================
% Bagian Konten Proposal Dimulai di Sini
% ======================================

% Mengatur ulang nomor halaman dan gaya halaman
\clearpage
\pagenumbering{arabic}
\setcounter{page}{1} 
\pagestyle{myfooter} % Ganti dari 'plain' ke 'myfooter'

% Atur margin standar (Opsional, jika Anda ingin margin isi berbeda dari sampul)
% \newgeometry{left=4cm, right=3cm, top=3cm, bottom=3cm} 
% Catatan: Perintah \newgeometry memerlukan paket 'geometry' yang sudah ada di cls.

\tableofcontents
\newpage
\section*{\docsumary}
Document Summary merupakan ringkasan yang hanya membahas isi dari dokumen capstone yang sedang dikerjakan, bukan keseluruhan proyek capstone. Bagian ini menjelaskan secara singkat tujuan, ruang lingkup, dan pokok pembahasan dari dokumen tersebut. Misalnya, pada dokumen EL1, ringkasan berfokus pada latar belakang, identifikasi masalah, dan lain-lain sesuai isi dari EL1.  Penulisan dilakukan secara padat dan jelas agar pembaca dapat langsung memahami inti dari dokumen tersebut tanpa mengaitkannya dengan bagian atau dokumen capstone lainnya.
\newpage



\section{\ProjectBenchmarkTitle}
Bagian ini menyajikan tinjauan komprehensif terhadap solusi, proyek, penelitian, atau paten terdahulu yang relevan dengan topik proyek tim Anda. Tujuannya adalah untuk memahami kondisi teknologi dan penelitian terkini, mengidentifikasi celah (gap) yang ada, dan menegaskan improvisasi dari proyek yang diusulkan.

\subsection{\PrevProjectsTitle}
Subbab ini membahas sistem atau proyek yang telah diimplementasikan, baik secara komersial maupun sebagai prototipe, yang memiliki fungsi serupa dengan proyek tim Anda. Fokus utama adalah pada fitur, kelebihan, kekurangan, dan metode implementasi yang digunakan oleh solusi-solusi tersebut. Analisis ini penting untuk mengidentifikasi area yang dapat dioptimalkan atau diperbaiki oleh proyek tim Anda.

\subsubsection*{Contoh Identifikasi Solusi Terdahulu:}
\begin{enumerate}
 
    \item \textbf{Sistem Pemantauan Kualitas Udara Komersial X (Misalnya: AirVisual Pro, uHoo):}
    \begin{enumerate}
        \item \textbf{Fitur:} Memantau CO$_2$, PM2.5, suhu, kelembaban. Terintegrasi dengan aplikasi mobile dan cloud. Beberapa memiliki rekomendasi ventilasi pasif.
        \item \textbf{Kelebihan:} Akurasi tinggi, data historis, user-friendly.
        \item \textbf{Kekurangan:} Harga mahal, aktuasi otomatis terbatas atau tidak ada (hanya rekomendasi), sensor gas spesifik (misalnya CH$_4$/LPG, CO) sering tidak tersedia dalam satu perangkat, sulit dikustomisasi.
        \item \textbf{Celah yang Diisi  ELSINTA:} \textbf{ ELSINTA} akan menawarkan sistem kontrol aktif (kipas/alarm) secara otomatis berdasarkan ambang batas gas spesifik (CO, CH$_4$) dengan biaya yang lebih terjangkau dan kemampuan kustomisasi yang tinggi untuk lingkungan spesifik.
    \end{enumerate}
    \item \textbf{Proyek Skripsi/Tugas Akhir Sejenis (Misalnya: "Sistem Deteksi Kebocoran Gas Berbasis Arduino"):}
    \begin{enumerate}
        \item \textbf{Fitur:} Menggunakan sensor MQ-2 untuk mendeteksi LPG dan mengaktifkan alarm buzzer. Tampilan sederhana di LCD.
        \item \textbf{Kelebihan:} Biaya rendah, mudah diimplementasikan.
        \item \textbf{Kekurangan:} Pemantauan terbatas pada satu jenis gas, tidak ada integrasi IoT/Cloud, tidak ada kontrol aktuator lain (misalnya kipas), tidak ada analisis data historis, akurasi sensor dasar seringkali rendah tanpa kalibrasi.
        \item \textbf{Celah yang Diisi  ELSINTA:} \textbf{ ELSINTA} mengintegrasikan multi-sensor gas (CO, CH$_4$, CO$_2$, PM2.5), fitur IoT untuk pemantauan real-time dan logging data, serta sistem kontrol multi-aktuator yang cerdas dan adaptif.
    \end{enumerate}
\end{enumerate}

\subsection{\RelResearchPatentsTitle}
Subbab ini mengulas penelitian ilmiah dan paten yang telah dipublikasikan terkait dengan teknologi atau metodologi yang akan digunakan dalam \textbf{ ELSINTA}. Analisis ini mencakup studi tentang efektivitas berbagai jenis sensor, algoritma kontrol otomatis, arsitektur sistem IoT untuk pemantauan lingkungan, atau metode kalibrasi sensor gas. Pemahaman terhadap penelitian dan paten ini membantu mengidentifikasi praktik terbaik (best practices) dan juga menunjukkan di mana \textbf{ ELSINTA} dapat memberikan kontribusi ilmiah atau inovatif.

\subsubsection*{Contoh Identifikasi Penelitian/Paten Relevan:}
\begin{enumerate}
    \item \textbf{Penelitian tentang Algoritma Kontrol PID untuk Ventilasi Otomatis:}
    \begin{enumerate}
        \item \textbf{Studi:} Penelitian oleh Smith et al. (2020) membahas optimalisasi kontroler PID untuk sistem ventilasi cerdas berdasarkan konsentrasi CO$_2$ dalam ruangan.
        \item \textbf{Relevansi dengan  ELSINTA:} \textbf{ ELSINTA} dapat mengadopsi atau memodifikasi algoritma serupa untuk mengendalikan kipas/ventilasi berdasarkan data dari multi-sensor (CO$_2$, CO, CH$_4$), dan membandingkan efektivitasnya. Hal ini membantu dalam perancangan modul kontrol aktuator yang responsif.
        \item \textbf{Kontribusi  ELSINTA:} Mengembangkan algoritma kontrol adaptif yang tidak hanya mempertimbangkan CO$_2$ tetapi juga prioritas bahaya dari gas lain seperti CO dan CH$_4$, serta kemungkinan kombinasi aktuator (kipas + alarm).
    \end{enumerate}
    \item \textbf{Paten tentang Metode Kalibrasi Sensor Gas Berbasis Machine Learning:}
    \begin{enumerate}
        \item \textbf{Studi/Paten:} Paten US1234567B2 (Inventor: Doe, 2021) menguraikan metode kalibrasi sensor gas semi-konduktor (MQ-series) menggunakan regresi adaptif untuk meningkatkan akurasi pembacaan pada berbagai kondisi lingkungan.
        \item \textbf{Relevansi dengan  ELSINTA:} Akurasi sensor MQ-series sering menjadi masalah. \textbf{ ELSINTA} dapat mengimplementasikan metode kalibrasi yang terinspirasi dari paten ini atau mengembangkan metode kalibrasi on-site yang lebih sederhana namun efektif untuk meningkatkan reliabilitas data.
        \item \textbf{Kontribusi  ELSINTA:} Mengeksplorasi implementasi kalibrasi sensor yang efisien secara komputasi pada mikrokontroler atau di sisi cloud untuk meningkatkan akurasi sistem secara keseluruhan tanpa memerlukan peralatan lab yang mahal.
    \end{enumerate}
    \item \textbf{Jurnal tentang Implementasi Arsitektur IoT untuk Smart Building:}
    \begin{enumerate}
        \item \textbf{Studi:} Artikel oleh Chen et al. (2019) menjelaskan arsitektur IoT yang skalabel menggunakan protokol MQTT dan platform cloud (misalnya AWS IoT) untuk pemantauan lingkungan di gedung pintar.
        \item \textbf{Relevansi dengan  ELSINTA:} \textbf{ ELSINTA} akan mengadopsi arsitektur IoT yang serupa untuk memastikan data sensor dapat diunggah, disimpan, dan divisualisasikan secara efisien melalui platform cloud.
        \item \textbf{Kontribusi  ELSINTA:} Mendesain arsitektur IoT yang optimal untuk lingkungan laboratorium/ruangan spesifik ITERA dengan mempertimbangkan batasan sumber daya dan skalabilitas untuk aplikasi di masa depan.
    \end{enumerate}
\end{enumerate}

\section{\GapAnalysisTitle}
Berdasarkan tolok ukur proyek pada bagian sebelumnya, subbab ini akan secara spesifik mengidentifikasi kesenjangan atau celah yang ada pada solusi/penelitian terdahulu. Analisis ini sangat penting untuk memvalidasi improvisasi yang dihasilkan proyek yang diusulkan, dengan menunjukkan bagaimana proyek ini akan mengisi celah-celah tersebut dan memberikan kontribusi yang berarti.

\subsection{\LimitationsTitle}
Solusi-solusi  yang ada saat ini, baik produk komersial maupun prototipe hasil penelitian sebelumnya, memiliki beberapa keterbatasan kunci yang belum sepenuhnya teratasi. Contoh Keterbatasan:
\begin{enumerate}
    \item \textbf{Kurangnya Integrasi Multi-Sensor dan Kontrol Aktif Terpadu:} Banyak solusi komersial hanya fokus pada beberapa parameter (misalnya CO$_2$ dan PM2.5) tanpa integrasi sensor gas berbahaya seperti CO dan CH$_4$ secara komprehensif. Jika ada, sistem kontrol otomatis (misalnya kipas) seringkali bersifat tunggal dan tidak adaptif terhadap berbagai jenis ancaman gas secara bersamaan. Prototipe penelitian seringkali juga terbatas pada satu jenis deteksi tanpa kontrol aktif.
    \item \textbf{Akurasi dan Keandalan Sensor Semi-Konduktor yang Rendah (tanpa Kalibrasi):} Sensor gas low-cost (seperti seri MQ) yang umum digunakan dalam proyek-proyek prototipe seringkali memiliki akurasi dan stabilitas yang kurang optimal. Tanpa proses kalibrasi atau koreksi yang memadai, data yang dihasilkan bisa tidak dapat diandalkan, terutama dalam jangka panjang atau di lingkungan yang berubah-ubah.
    \item \textbf{Biaya Implementasi yang Tinggi untuk Fitur Lengkap:} Produk komersial yang menawarkan fitur lengkap (multi-sensor, IoT, data historis) cenderung memiliki harga yang sangat mahal, membuatnya tidak terjangkau untuk aplikasi yang lebih luas atau spesifik. Di sisi lain, solusi low-cost seringkali mengorbankan fungsionalitas dan keandalan.
    \item \textbf{Keterbatasan Kustomisasi dan Adaptasi Lingkungan:} Solusi yang ada seringkali bersifat "plug-and-play" dengan konfigurasi terbatas. Ini menyulitkan adaptasi sistem untuk kebutuhan spesifik lingkungan ruangan (misalnya, perbedaan ukuran ruangan, jenis polutan dominan, atau ambang batas yang disesuaikan).
    \item \textbf{Kurangnya Visualisasi Data yang Informatif dan Akses Real-time:} Beberapa proyek terdahulu mungkin memiliki fitur pemantauan, tetapi seringkali terbatas pada tampilan lokal (LCD) tanpa kemampuan visualisasi data real-time melalui dashboard berbasis web/mobile atau penyimpanan data historis untuk analisis jangka panjang.
\end{enumerate}

\subsection{\ProjectContribsTitle}
Sub-bab ini secara eksplisit mengidentifikasi dan mendokumentasikan \textbf{kontribusi} yang dihasilkan oleh Proyek Capstone ini. Contoh celah yang diisi oleh project:
\begin{enumerate}
    \item \textbf{Sistem Multi-Sensor Terintegrasi dengan Kontrol Aktif Adaptif:} \textbf{Project ELSINTA} akan mengintegrasikan sensor untuk CO, CH$_4$/LPG, CO$_2$, dan PM2.5 dalam satu unit, dilengkapi dengan algoritma kontrol cerdas yang secara otomatis mengaktifkan berbagai aktuator (kipas/ventilasi, alarm suara/visual) berdasarkan tingkat konsentrasi masing-masing gas dan prioritas bahayanya. Ini berbeda dari sistem tunggal atau pasif yang ada.
    \item \textbf{Peningkatan Akurasi Melalui Kalibrasi dan Koreksi Data:} Proyek ini akan mengimplementasikan metode kalibrasi awal dan/atau algoritma koreksi data sederhana pada firmware mikrokontroler atau di cloud untuk meningkatkan akurasi dan stabilitas pembacaan sensor MQ-series, membuatnya lebih andal dibandingkan prototipe tanpa kalibrasi.
    \item \textbf{Solusi Efektif Biaya dengan Fungsionalitas Tinggi:} \textbf{Project ELSINTA} bertujuan untuk membangun prototipe yang menawarkan fungsionalitas komprehensif (pemantauan multi-gas, IoT, kontrol otomatis) dengan biaya implementasi yang jauh lebih rendah dibandingkan produk komersial sejenis, sehingga lebih mudah diakses untuk implementasi di lingkungan pendidikan atau UKM.
    \item \textbf{Desain Moduler dan Fleksibilitas Konfigurasi:} Sistem akan dirancang dengan pendekatan modular yang memungkinkan kustomisasi sensor dan aktuator, serta konfigurasi ambang batas alarm dan logika kontrol yang dapat disesuaikan dengan kebutuhan spesifik lingkungan ruangan yang berbeda.
    \item \textbf{Platform Monitoring Real-time dengan Analisis Data Historis:} \textbf{Project ELSINTA} akan menyediakan dashboard berbasis web/mobile yang intuitif untuk visualisasi data kualitas udara secara real-time dan penyimpanan data historis di cloud. Ini memungkinkan pengguna untuk melakukan analisis tren jangka panjang dan mengambil keputusan yang lebih baik terkait manajemen kualitas udara.
\end{enumerate}



\section{\StandardsComplianceTitle}

Bagian ini menguraikan standar teknis, regulasi, dan aturan yang harus dipenuhi oleh \textbf{Project ELSINTA (Sistem Gas Quality Monitoring and Controlling Dalam Ruangan)}. Mematuhi standar menjamin kualitas, keamanan, dan interoperabilitas sistem, sementara mematuhi regulasi memastikan legalitas dan kepatuhan terhadap keselamatan lingkungan dan kerja.

\subsection{\TechnicalStandardsTitle}
\textit{Penjelasan: Subbab ini harus mengidentifikasi dan menjelaskan standar teknis yang relevan dengan aspek-aspek kunci Project ELSINTA, seperti komunikasi \text{IoT}, sensor, modul nirkabel, atau metodologi pengembangan perangkat lunak. Standar ini berfungsi sebagai pedoman desain dan pengujian untuk memastikan kualitas, interoperabilitas, dan performa sistem.}

Standar teknis adalah pedoman yang diadopsi secara luas untuk memastikan konsistensi, keandalan, dan interoperabilitas komponen sistem.

\subsubsection*{Contoh Standar Teknis yang Relevan:}
\begin{enumerate}
    \item \textbf{\text{IEEE 802.11} (WiFi):} Standar ini wajib dipenuhi untuk modul komunikasi nirkabel ($\text{WiFi}$) yang digunakan oleh mikrokontroler ($\text{NodeMCU/ESP32}$) dalam Project ELSINTA. Kepatuhan terhadap $\text{IEEE 802.11}$ memastikan perangkat dapat terhubung dengan jaringan internet dan mengirimkan data secara standar ke \emph{cloud}.
    \item \textbf{\text{ISO 9001} (Sistem Manajemen Mutu):} Meskipun Project ELSINTA adalah prototipe, prinsip-prinsip sistem manajemen mutu $\text{ISO 9001}$ dapat digunakan sebagai pedoman dalam proses perancangan, pengembangan, dokumentasi, dan pengujian sistem. Hal ini bertujuan untuk memastikan kualitas luaran (\emph{deliverables}) dan metodologi kerja yang terstruktur.
    \item \textbf{\text{IEEE 1451} (Transducer Interface Standards):} Standar ini relevan untuk mendefinisikan \emph{interface} cerdas antara berbagai sensor (misalnya $\text{CO}, \text{CO}_2, \text{PM2.5}$) dan mikrokontroler. Penerapan prinsip $\text{IEEE 1451}$ dapat memfasilitasi pertukaran data yang efisien, kalibrasi sensor, dan modularitas sistem sensor.
    \item \textbf{\text{MQTT} (Message Queuing Telemetry Transport):} Sebagai protokol standar \emph{de facto} yang ringan untuk komunikasi \text{IoT}, \text{MQTT} akan digunakan dalam Project ELSINTA untuk pengiriman data sensor yang efisien dan real-time dari perangkat ke broker di \emph{cloud}. Protokol ini dipilih karena efisiensinya dalam penggunaan bandwidth dan sumber daya.
\end{enumerate}

\subsection{\RegulationsTitle}
\textit{Penjelasan: Subbab ini harus mengidentifikasi regulasi atau aturan hukum yang wajib dipatuhi oleh Project ELSINTA. Fokus pada regulasi yang berkaitan dengan keselamatan dan kesehatan kerja, kualitas udara, serta penggunaan perangkat telekomunikasi di Indonesia. Ini menunjukkan bahwa proyek Anda tidak hanya berfungsi, tetapi juga aman dan legal.}

Regulasi adalah aturan hukum yang wajib dipatuhi. Dalam konteks proyek pemantauan kualitas udara dan keselamatan, regulasi ini terkait erat dengan keselamatan dan kesehatan kerja serta standar lingkungan.

\subsubsection*{Contoh Regulasi yang Relevan:}
\begin{enumerate}
    \item \textbf{\text{Permenaker No. 5 Tahun 2018} (Keselamatan dan Kesehatan Kerja Lingkungan Kerja):} Regulasi ini mengatur Nilai Ambang Batas (\text{NAB}) paparan zat kimia di lingkungan kerja, termasuk $\text{CO}$ dan $\text{CO}_2$. Project ELSINTA harus mengacu pada regulasi ini untuk menetapkan ambang batas alarm $\text{CO}$ (misalnya 50 ppm selama 8 jam kerja) dan gas lainnya guna memastikan lingkungan kerja yang aman.
    \item \textbf{Regulasi \text{SNI} tentang Kualitas Udara Dalam Ruangan (misalnya \text{SNI 7456}):} Mengatur standar umum untuk kualitas udara di berbagai jenis bangunan. Project ELSINTA akan menggunakan standar \text{SNI} yang berlaku untuk menetapkan batas aman $\text{CO}_2$ (biasanya di bawah 1000 ppm) demi kenyamanan dan kesehatan penghuni ruangan.
    \item \textbf{Aturan \text{SDPPI} (Direktorat Jenderal Sumber Daya dan Perangkat Pos dan Informatika):} Jika komponen perangkat nirkabel ($\text{WiFi}$ atau teknologi nirkabel lain) yang digunakan dalam Project ELSINTA adalah impor, perangkat tersebut harus mematuhi aturan \text{SDPPI} terkait sertifikasi dan penggunaan spektrum frekuensi radio di Indonesia.
    \item \textbf{Aturan Institusi ITERA (Keselamatan Laboratorium):} Project ELSINTA harus mematuhi semua \text{SOP} (Standard Operating Procedure) dan aturan \text{K3} (Keselamatan dan Kesehatan Kerja) yang ditetapkan oleh ITERA, terutama jika sistem akan diimplementasikan atau diuji di lingkungan laboratorium atau fasilitas ITERA.
\end{enumerate}

\section{\ReferencesTitle}
Daftar pustaka ditulis sesuai standar IEEE, sebaiknya gunakan Citation tools seperti Mendeley, Zotero dll.  
Contoh format IEEE:  
\renewcommand{\refname}{}      % Hapus judul default "Pustaka"
\bibliographystyle{IEEEtran} % Menggunakan style sitasi IEEE
\bibliography{pustaka}    % Memanggil file pustaka.bib





\section{\AbbrevTitle}

% longtable tidak menggunakan lingkungan table[] di sekitarnya.
% Perintah \begin{longtable} langsung menggantikan \begin{table}\begin{tabular}
\begin{longtable}{|p{4cm}|p{11cm}|} % Lebar kolom ditetapkan agar teks wrap
    \caption{Daftar Singkatan (Abbreviations)} \label{tab:abbreviations2} \\
    \hline
    \textbf{Abbreviation} & \textbf{Definition} \\
    \hline
    \endhead % Mengakhiri baris header yang akan diulang di setiap halaman
    
    \hline 
    \multicolumn{2}{|r|}{Lanjutan di halaman berikutnya...} \\ 
    \endfoot % Footer untuk semua halaman kecuali yang terakhir
    
    \hline
    \endlastfoot % Footer untuk halaman terakhir (kosong)

    % --- Isi Tabel ---
    IoT & Internet of Things \\
    \hline
    KPI & Key Performance Indicator \\
    \hline
    ROI & Return on Investment \\
    \hline
    WBS & Work Breakdown Structure \\
    \hline
    SWOT & Strengths, Weaknesses, Opportunities, Threats \\
    \hline
    BMC & Business Model Canvas \\
    \hline
    IEEE & Institute of Electrical and Electronics Engineers \\
    \hline
    ISO & International Organization for Standardization \\
    \hline
    G-MoCo & Gas Quality Monitoring and Controlling (Nama Proyek Anda) \\
    \hline
    CO & Carbon Monoxide \\
    \hline
    CH$_4$ & Methane (atau LPG) \\
    \hline
    PM2.5 & Particulate Matter 2.5 \\
    \hline
    
    % Tambahkan lebih banyak baris di sini untuk memicu pemisahan halaman
    % ... (Tambahkan entri lain jika tabel perlu memanjang)

\end{longtable}

% Tambahkan \newpage jika Anda ingin halaman isi dimulai di halaman berikutnya
% \newpage
% \section*{Pendahuluan}
% Di sinilah konten proposal Anda akan dimulai...

\end{document}